\begin{claim}{C.66}{Proved}{Zero-dimensional frontiers give a global centralizer filling}
Let $f_t$ be a based ambient-null loop in
$\operatorname{Homeo}^+(\R^2)$, let $G$ be the group
generated by its values, and let $H$ be any group
of plane homeomorphisms commuting with every $f_t$.
Let $Y$ be the closed union of the nonopen $G$-orbits.
When $Y$ is zero-dimensional, the original loop
has a based disk filling in
$C_{\operatorname{Homeo}(\R^2)}(H)$.
Every map of the filling fixes $Y$ pointwise.
This includes the empty frontier and imposes no
discreteness or local compactness assumption on $H$.

The topological comparison used here is the following
actual isomorphism. For every nonempty compact
zero-dimensional set $Z\subset S^2$, put $M=S^2\setminus Z$.
With compact-open group topologies, including inverse
convergence, restriction is an isomorphism of topological
groups
\[
 R_Z:\operatorname{Homeo}(S^2;Z)
             \xrightarrow{\ \cong\ }\operatorname{Homeo}(M),
 \qquad \operatorname{Homeo}(S^2;Z)
       =\{a:a(Z)=Z\}.                                  \tag{C.66.1}
\]
Its inverse is the unique end extension.
A connected parameter family containing the identity
extends to a family fixing every point of $Z$.

Consequently, for a connected admissible surface of C.1,
an ambient-null loop in the centralizer of an arbitrary
specified operator action is null-homotopic in that
same centralizer whenever the nonopen locus of its
based plane lift is zero-dimensional. All original
operator and deck labels are retained.
\end{claim}

\begin{proof}
\emph{Finite disjoint Jordan covers of the deleted set.}
We first prove the planar compactification statement
with all the topology needed for parameter families.
End extension for zero-dimensional deleted sets is
also used in \cite[Proposition~1.5 and Appendix~A.1]
{C_BeguinCrovisierLeRoux2020}; the following construction
supplies the actual group topology comparison.

Let $Z\subset S^2$ be compact and zero-dimensional.
It is a proper subset. Choose a stereographic pole
outside $Z$, so $Z$ becomes a compact subset of the
coordinate plane. Given an open neighborhood $V$ of
$Z$ and $\varepsilon>0$, there is a finite family
of closed Jordan disks $D_1,\ldots,D_r$ such that
\[
 \begin{gathered}
 D_i\cap D_j=\varnothing\quad(i\ne j),\qquad
 Z\subset\bigcup_j\operatorname{int}D_j,\qquad
 D_j\subset V,\qquad D_j\cap Z\ne\varnothing,\\
              \operatorname{diam}_{S^2}D_j<\varepsilon .
 \end{gathered}                                         \tag{C.66.2}
\]
Here and below a diameter refers to one fixed compatible
metric on the compact sphere.

To construct the disks, partition $Z$ into finitely
many disjoint compact relatively clopen sets of
arbitrarily small coordinate diameter. Such a partition
exists by the clopen basis and compactness: take a finite
clopen cover of the required mesh and disjointize it by
successive differences. Distinct nonempty pieces have
positive distance. Cover them by a finite sufficiently
fine collection of coordinate grid cells and take a
small polygonal regular neighborhood of that union.
One obtains the neighborhood by separating corner
contacts in tiny vertex neighborhoods and rounding
the remaining boundary corners. Choose the cell mesh
and rounding size below the separation of the clopen
pieces. Each connected component of this regular
neighborhood then lies in a small neighborhood of one
partition piece. Delete any component not meeting $Z$.
Its other components are compact connected planar
surfaces, with finitely many disjoint Jordan boundary
curves, and they still contain $Z$ in their interiors.

Fill each component up to its outer Jordan curve.
The resulting disk lies in the convex hull of that
component: a point outside the convex hull is separated
by a line and belongs to the unbounded component of
the complement of its outer curve. Its coordinate
diameter is therefore bounded by that of the original
component. The outer curves of different components
are disjoint. Their closed Jordan disks are either
disjoint or nested; keep only the maximal disks.
Their closures are pairwise disjoint, and they cover
all the original components and hence $Z$.
The partition mesh and neighborhood sizes were arbitrary.
Choose them small enough that each retained disk,
which contains a point of $Z$, lies in $V$ and has
spherical diameter below $\varepsilon$. This proves
(C.66.2), without a local-connectivity assumption on $Z$.

The complement of $Z$ is path connected. For two
points outside $Z$, apply (C.66.2) with a neighborhood
avoiding both points. The complement of the interiors
of the finitely many disjoint disks is a connected
sphere with finitely many disks removed, and it
contains both points. A path in this complement
avoids $Z$. The same argument inside a Jordan disk
$D$ with $\partial D\cap Z=\varnothing$ shows that
$\operatorname{int}D\setminus Z$ is path connected:
cover the compact set $Z\cap D$ by finitely many
smaller disks in its interior, avoiding the two
chosen points, and join them in the remaining
finitely connected disk.

\emph{The exact end compactification.}
Apply (C.66.2) successively to obtain disk families
$D_{n,j}$ of mesh less than $2^{-n}$, with every
disk at stage $n+1$ contained in the interior of
one disk at stage $n$. Every disk meets $Z$.
Set
\[
 K_n=S^2\setminus\bigcup_j\operatorname{int}D_{n,j},
 \qquad
 U_{n,j}=\operatorname{int}D_{n,j}\setminus Z .
                                                               \tag{C.66.3}
\]
The sets $K_n$ are compact connected planar surfaces
contained in $M$, satisfy $K_n\subset\operatorname{int}K_{n+1}$,
and exhaust $M$. Indeed each disk meets $Z$ and its
diameter tends to zero, so every point outside $Z$
eventually lies outside all the disks. The preceding
connectivity argument shows that the $U_{n,j}$ are
exactly the components of $M\setminus K_n$.
Each is non-relatively-compact in $M$, since its
disk meets $Z$.

An end is a nested compatible choice of these
components. The corresponding nested closed disks
have diameters tending to zero. Their intersection
is one point of $Z$: each contains a point of $Z$,
and compactness and nesting give an intersection
with $Z$. Conversely every point of $Z$ determines
one such choice. Thus the map
\[
 \Xi:\widehat M\longrightarrow S^2,\qquad
 \Xi|_M=\id,\qquad
 \Xi((U_{n,j(n)})_n)=\bigcap_nD_{n,j(n)}
                                                               \tag{C.66.4}
\]
is a bijection from the end compactification.
A basic end neighborhood consisting of $U_{n,j}$
and the ends it contains maps to
$\operatorname{int}D_{n,j}$. These disks form a
neighborhood basis at their limiting point, by
their vanishing diameters. Equation (C.66.4)
and its inverse are therefore continuous,
and the topology on $M$ is unchanged.

Every $a\in\operatorname{Homeo}(M)$ is proper.
Given a target $K_m$, the compact set $a^{-1}(K_m)$
lies inside some $K_n$. The connected set
$a(U_{n,j})$ then lies in a unique component
of $M\setminus K_m$. Following a nested sequence
of source components defines its image end.
The same construction for $a^{-1}$ gives the
inverse map, and the basic neighborhoods just
described prove continuity. This is the unique
extension $\widehat a$ to $S^2$: $M$ is dense
because a zero-dimensional subset has empty
interior. The component construction, or
uniqueness on this dense set, gives
$\widehat{ab}=\widehat a\,\widehat b$ and
$\widehat{a^{-1}}=(\widehat a)^{-1}$.
The extension preserves $Z$ setwise.
Conversely a sphere homeomorphism preserving
$Z$ restricts to $M$ and is its unique extension.
This proves the algebraic bijection in (C.66.1).

\emph{The compact-open and inverse topology.}
Restriction in (C.66.1) is continuous on compact
subsets of $M$, for maps and inverses. We prove
continuity of extension at the identity; the
group homomorphism then gives continuity everywhere.
Both groups are metrizable, so it suffices to take
a sequence $a_i\to\id$ in the actual homeomorphism
group of $M$, including $a_i^{-1}\to\id$ on compact sets.

Fix $\varepsilon>0$ and choose $m$ so all disks
$D_{m,j}$ have diameter below $\varepsilon$.
The set
\[
             K_m\ \cup\ \bigcup_i a_i^{-1}(K_m)
                                                               \tag{C.66.5}
\]
is compact in $M$. To see this exactly, the
convergent sequence $\{\id,a_i^{-1}\}$ is a compact
parameter space, and its evaluation on the compact
set $K_m$ is a compact image. Choose $N\ge m$
with (C.66.5) contained in $K_N$.
Every $a_i(U_{N,j})$ consequently avoids $K_m$.
It is connected and hence lies in one component
of $M\setminus K_m$.

Choose a point $p_j$ in each of the finitely many
$U_{N,j}$. This set lies inside one parent
$U_{m,\ell(j)}$. Since $a_i(p_j)\to p_j$,
for all sufficiently large $i$ the component
containing $a_i(U_{N,j})$ is exactly that parent.
The end-extension construction then also sends
every end in $D_{N,j}$ into $D_{m,\ell(j)}$.
Thus each point of $\operatorname{int}D_{N,j}$,
including its deleted points, and its image
under $\widehat a_i$ lie in the same disk
$D_{m,\ell(j)}$ of diameter below $\varepsilon$.
On the remaining compact set $K_N$, the
original maps $a_i$ converge uniformly to identity.
It follows that $\widehat a_i\to\id$ uniformly
on the whole sphere. Apply the identical argument
to the sequence $a_i^{-1}$ to obtain uniform
inverse convergence. This proves (C.66.1) as
an isomorphism of the stated topological groups.

In particular every continuous parameter family
extends jointly with its inverse. When its parameter
space is connected and contains a parameter at which
the map is identity, evaluation at any $z\in Z$
is a continuous map into the zero-dimensional
space $Z$ with connected image. That image is
the singleton $\{z\}$. The entire family fixes
every deleted point individually. This is the
pointwise conclusion used below; the full group
in (C.66.1) preserves $Z$ only setwise.

\emph{Application to the original plane loop.}
The group $G$ is path connected: any finite word
in loop values is joined to identity by scaling
each of its time parameters to zero. Its orbits
are therefore path connected. A nonopen orbit
contained in the zero-dimensional set $Y$
must be a point, so $G$ fixes $Y$ pointwise.
Let
\[
                  Z=Y\cup\{\infty\}\subset S^2 .
                                                               \tag{C.66.6}
\]
This is compact and zero-dimensional. For
completeness, a closed zero-dimensional subset
of the plane is locally compact and has a basis
of compact relatively open-and-closed sets:
inside a compact neighborhood choose a clopen
basis element with closure still in that
neighborhood. Finite unions of such sets give
an increasing compact clopen exhaustion of $Y$.
Their complements in $Z$ form a clopen basis
at the point at infinity. The old clopen basis
at each finite point is unchanged. The empty
$Y$ gives the singleton $Z=\{\infty\}$.

The complement $O=\R^2\setminus Y=S^2\setminus Z$
is connected by the disk-cover argument.
It is a union of disjoint open $G$-orbits.
Connectedness makes it one such orbit.
Every $h\in H$ transports $G$-orbits to $G$-orbits
with their actual labels, since it commutes
with $G$. It preserves $Y$ and the unique open
orbit $O$.

Let $P$ be the actual restriction image of
$H\to\operatorname{Homeo}(O)$, retaining its
kernel $\{h:h|_O=1\}$, and let $Q$ be the
closure of $P$ in that homeomorphism group.
C.30 makes the full local centralizer of
$G|_O$ a closed locally compact group acting
freely and properly on $O$. The closed subgroup
$Q$ of this full centralizer has these same
properties. The restricted original loop
commutes with $Q$.

C.32 proves that this restriction is null-homotopic
in $\operatorname{Homeo}(O)$, using the specified
ambient-null input. C.33 therefore gives an actual
based disk filling
\[
 A:D^2\longrightarrow
          C_{\operatorname{Homeo}(O)}(Q),
       \qquad A|_{\partial D^2}=f_t|_O .
                                                               \tag{C.66.7}
\]
Apply the inverse of (C.66.1) to this whole family.
The disk is connected and contains the based
identity parameter, so every extended map fixes
all of $Z$, including the specified point
$\infty$. Restricting these sphere maps to the
original plane gives a jointly continuous family
of plane homeomorphisms and inverses, fixing $Y$.
On $O$ it commutes with each restriction of
each original $h\in H$. The two compositions
are continuous on the plane and agree on the
dense set $O$, so they agree everywhere.
The boundary family is the original $f_t$ by
the same density argument. This proves the
claimed filling with every operator label retained.

\emph{The actual lift and descent for an admissible surface.}
Let $p:\R^2\to U$ be the universal covering of a
connected admissible surface, and start with a
based centralizer loop which is null-homotopic
in its ambient homeomorphism group.
Lift that ambient disk family after composition
with $p$, choosing the identity at its based
parameter. The covering lifting criterion and
the corresponding inverse family give actual
plane homeomorphisms. Its boundary is a closed
ambient-null based plane loop $\widetilde f_t$.

Take $H$ to consist of every original lift of
every specified surface operator, including
all deck transformations. For each such lift
$h$, the maps $\widetilde f_t h$ and
$h\widetilde f_t$ lift the same surface family
and agree at the based parameter; uniqueness
of the covering lift makes them equal for
every $t$. Under the zero-dimensional locus
assumption the plane result supplies a disk
family $K_w$ commuting with every one of
these original lifts. Commutation with the
deck transformations gives the unique descent
\[
                    \overline K_w\,p=pK_w .
                                                               \tag{C.66.8}
\]
The open covering quotient makes the descended
map and inverse families jointly continuous.
Commutation with each lifted operator descends
to commutation with its specified surface
operator. The family has the original boundary.

For $U=X\setminus S$ with nonempty finite $S$,
the same end-extension topology argument, now
using disjoint coordinate disks about the
finitely many punctures in $X$, extends this
connected based family to $X$ fixing $S$
pointwise. In the compact case there is no
extension step. Thus (C.66.8) is a disk filling
in the original group $A^0\cap A^u$, with
its original operator action. No replacement
of that group by an abstract quotient is made.

The proof handles zero-dimensional $Y$ through
its exact end compactification. It does not
extend an arbitrary intrinsic filling across
a nontrivial line or circle frontier, or across
an arbitrary positive-dimensional fixed set.
Those remaining extensions in C.12 are unchanged.
\end{proof}
